% =====================================================================
%  School Wits - Cambridge O Level Additional Mathematics 4037
%  Paper 2, 4037/22/M/J/16
%
%  Keep 4037_s16_qp_22.pdf and 3-01.png beside this file (upload both to Overleaf). Nothing else is
%  needed: the diagrams, tables and stacked fractions are clipped out of the
%  official PDF by pdflatex itself, at the exact size and position they have
%  in the original.
%
%  Build modes - uncomment one line inside the document:
%     \swwbtrue     the printed workbook head
% =====================================================================
\documentclass[11pt,a4paper]{article}

\usepackage[paperwidth=595.276bp, paperheight=841.89bp,
            left=49.6bp, right=46.376bp,
            top=62.56bp, textheight=734bp,
            head=42.4bp, headsep=7bp, footskip=22bp]{geometry}

\usepackage[T1]{fontenc}
\usepackage[utf8]{inputenc}
\renewcommand{\rmdefault}{ptm}\renewcommand{\familydefault}{\rmdefault}
% The paper is set in Times.  Nimbus Roman has Times's advances exactly; its
% ligatures and pair kerning are switched off, as the paper has none.
\usepackage[protrusion=false,expansion=false,kerning=false,spacing=false,tracking=false]{microtype}
\DisableLigatures{encoding = *, family = *}
\usepackage{amsmath}\usepackage{amssymb}
\usepackage{graphicx}
\usepackage{xcolor}
\usepackage{tikz}
\usepackage{fancyhdr}
\usepackage{calc}
\definecolor{swrule}{gray}{0.62}

% The paper sets its body at 11.5bp on a 13bp baseline.  The class would reset
% both at \begin{document}, so \normalsize itself is pinned to the measurement.
\makeatletter
\renewcommand\normalsize{\@setfontsize\normalsize{11.5bp}{13bp}%
  \abovedisplayskip\z@ \belowdisplayskip\z@
  \abovedisplayshortskip\z@ \belowdisplayshortskip\z@}
\makeatother
\setlength{\parindent}{0pt}
\setlength{\parskip}{0pt}
\linespread{1}\normalbaselineskip=13bp \baselineskip=13bp
\lineskiplimit=-\maxdimen \lineskip=0pt
\makeatletter\def\@listi{\leftmargin\leftmargini}\makeatother
\setlength{\topskip}{10.04bp}
% A page that opens lower than the standard first baseline needs its offset
% as non-discardable space: TeX throws away ordinary glue at the top.

\clubpenalty=10000 \widowpenalty=10000
\hyphenpenalty=10000 \exhyphenpenalty=10000 \tolerance=4000 \emergencystretch=0pt
%  The paper's word space, measured: 2.875bp, stretching by up to
%  2.85bp and shrinking by up to 1.2bp to justify a line.
\frenchspacing
\spaceskip=2.875bp plus 2.85bp minus 1.2bp
\xspaceskip=2.875bp plus 2.85bp minus 1.2bp
\hbadness=10000
% a line the paper fills to a hair past the measure is not worth a warning
\hfuzz=1bp

% ---- symbols, set the way the paper sets them ------------------------
\newcommand{\ftsize}{\fontsize{8bp}{9bp}\selectfont}
\newcommand{\ohms}{{\usefont{U}{psy}{m}{n}W}}
\newcommand{\dg}{\textdegree}
\newcommand{\degC}{\textdegree C}
\newcommand{\x}{\ensuremath{\times}}
% super- and subscripts as the paper sets them: 8.05bp, 5.2bp up or 4.0bp
% down, taking no height so a line holding one keeps its place
\newcommand{\pow}[1]{\raisebox{5.2bp}[0pt][0pt]{\fontsize{8.05bp}{8.05bp}\selectfont #1}}
\newcommand{\sub}[1]{\raisebox{-4.0bp}[0pt][0pt]{\fontsize{8.05bp}{8.05bp}\selectfont #1}}
%  a nuclide symbol: mass number above, proton number below, as the paper sets it
% a nuclide symbol: the two numbers centred over one another.  The paper's
% own offsets and size come in the optional argument, up/down/size in bp,
% because they differ from paper to paper: \stack[3.86/3.84/8.25]{X}{Y}
\makeatletter
\newcommand{\stack}[3][5.34/2.66/8]{\sw@stack#1\relax{#2}{#3}}
\def\sw@stack#1/#2/#3\relax#4#5{\begingroup\fontsize{#3bp}{#3bp}\selectfont
  \setbox0=\hbox{#4}\setbox2=\hbox{#5}%
  \dimen0=\ifdim\wd0>\wd2 \wd0\else\wd2\fi
  \rlap{\raisebox{#1bp}[0pt][0pt]{\hbox to\dimen0{\hss\box0\hss}}}%
  \raisebox{-#2bp}[0pt][0pt]{\hbox to\dimen0{\hss\box2\hss}}\endgroup}
\makeatother
% a few characters the paper sets in another font - Times New Roman for Roman
% numerals - set in that font's free clone at the paper's size: \swf{ptm}{12.00}{VII}
\newcommand{\swf}[3]{{\fontfamily{#1}\fontsize{#2bp}{#2bp}\selectfont #3}}
%  the paper's thin space inside units, measured: 1.44bp (LaTeX's own is 1.83bp)
\AtBeginDocument{\renewcommand{\,}{\ifmmode\mskip\thinmuskip\else\kern1.44bp\relax\fi}}

% =====================================================================
%  TOPIC COLOURS — one definition, read by the chip and by the website
% =====================================================================
\definecolor{swDefault}{HTML}{D8DDE4}
\definecolor{tp403701}{HTML}{D44949}
\definecolor{tp403702}{HTML}{49A6D4}
\definecolor{tp403703}{HTML}{D4C249}
\definecolor{tp403704}{HTML}{9A49D4}
\definecolor{tp403705}{HTML}{49D449}
\definecolor{tp403706}{HTML}{D4498F}
\definecolor{tp403707}{HTML}{D48F49}
\definecolor{tp403708}{HTML}{4978D4}
\definecolor{tp403709}{HTML}{8FD449}
\definecolor{tp403710}{HTML}{D449D4}
\definecolor{tp403711}{HTML}{49D4A6}
\definecolor{tp403712}{HTML}{D46C49}
\definecolor{tp403713}{HTML}{6049D4}
\definecolor{tp403714}{HTML}{C9D449}
\definecolor{tp403715}{HTML}{49C9D4}
\definecolor{tp403716}{HTML}{D4496C}
\definecolor{tp403717}{HTML}{49D478}
\definecolor{tp403718}{HTML}{D4AB49}
\definecolor{tp403719}{HTML}{4955D4}

\makeatletter
\newcommand{\sw@key}[1]{sw@tc@\expandafter\zap@space\detokenize{#1} \@empty}
\newcommand{\deftopic}[2]{%
  \protected@edef\sw@tmp{#1}%
  \expandafter\sw@deftopic\expandafter{\sw@tmp}{#2}}
\newcommand{\sw@deftopic}[2]{%
  \expandafter\xdef\csname\sw@key{#1}\endcsname{#2}}
\makeatother

\deftopic{Functions}{tp403701}
\deftopic{Quadratic Functions}{tp403702}
\deftopic{Factors of Polynomials}{tp403703}
\deftopic{Equations, Inequalities \& Graphs}{tp403704}
\deftopic{Simultaneous Equations}{tp403705}
\deftopic{Logarithmic \& Exponential Functions}{tp403706}
\deftopic{Straight-Line Graphs}{tp403707}
\deftopic{Coordinate Geometry of the Circle}{tp403708}
\deftopic{Circular Measure}{tp403709}
\deftopic{Trigonometry}{tp403710}
\deftopic{Permutations \& Combinations}{tp403711}
\deftopic{Binomial Theorem}{tp403712}
\deftopic{Arithmetic \& Geometric Progressions}{tp403713}
\deftopic{Vectors in Two Dimensions}{tp403714}
\deftopic{Differentiation}{tp403715}
\deftopic{Integration}{tp403716}
\deftopic{Calculus for Kinematics}{tp403717}
\deftopic{Indices \& Surds}{tp403718}
\deftopic{Excluded from the present curriculum}{tp403719}

% =====================================================================
%  THE HEADER — the ONE thing that differs between the two builds
% =====================================================================
\makeatletter
\def\sw@firstof#1\textperiodcentered#2\@nil{#1}
\newcommand{\swfirsttopic}[1]{\sw@firstof#1\textperiodcentered\@nil}
% every topic of a question as its own chip, main topic first, 4bp apart:
% the list is split where \examq's topics are joined, at "\textperiodcentered\ "
\newcommand{\topicchips}[1]{\sw@chipsloop#1\textperiodcentered\ \@nil}
\def\sw@chipsloop#1\textperiodcentered\ #2\@nil{\topicchip{#1}%
  \if\relax\detokenize{#2}\relax\else\hspace{4bp}\sw@chipsloop#2\@nil\fi}
\newcommand{\sw@hstrut}{\rule[-3.4bp]{0pt}{14.6bp}}
\newcommand{\topicchip}[1]{%
  \protected@edef\sw@tmp{#1}%
  \expandafter\sw@chip\expandafter{\sw@tmp}}
%  The chip: a rounded tag in three shades of the topic's own colour - a
%  tinted fill, a hairline border, and text darkened into the same hue, so it
%  reads as a label rather than a block of paint.  The box keeps the old
%  measurements exactly: 9bp above the baseline, 2.9bp below, 4.9bp of air
%  either side of the text, so nothing around it moves.
%
%  Three knobs, if the look wants adjusting:
%    sw@bg  fill      !60!white is a soft tint; raise towards !100 for solid
%    sw@bd  border    a lightened version of the text colour
%    sw@ink text      !45!black keeps it legible on any of the pastels
%  For a solid chip instead, set sw@bg to sw@base and drop the draw= key.
\newcommand{\sw@chip}[1]{%
  \begingroup
  \edef\sw@col{\ifcsname\sw@key{#1}\endcsname
      \csname\sw@key{#1}\endcsname\else swDefault\fi}%
  \colorlet{sw@base}{\sw@col}%
  \colorlet{sw@ink}{sw@base!45!black}%
  \colorlet{sw@bg}{sw@base!60!white}%
  \colorlet{sw@bd}{sw@ink!35!white}%
  \begin{tikzpicture}[baseline=(swchip.base)]
    \node[anchor=base, at={(0,0)},
          inner xsep=4.9bp, inner ysep=0bp,
          text height=9bp, text depth=2.9bp,
          rounded corners=2.6bp,
          fill=sw@bg, draw=sw@bd, line width=0.4bp,
          text=sw@ink,
          font=\fontsize{8.6bp}{8.6bp}\selectfont\bfseries] (swchip) {#1\unskip};
  \end{tikzpicture}%
  \endgroup}

% =====================================================================
%  THE QUESTION HEAD
%  Every question carries \examq{<paper id>}{<number>}{<topics>}{<marks>}.
%  It is the single source of truth: the website splits the file on it and
%  filters on the topic list, the contents page groups by the first topic,
%  and \total reads the marks from it.
%
%  Three builds, one file:
%     (default)    the yearly paper - the head prints NOTHING, so the page
%                  stays a copy of the official one, page for page
%     \swwbtrue    the printed workbook - paper id, question number and every
%                  topic as a colour chip, main topic first; no marks
%     \swwebtrue   the website - the full line, topics and marks included
% =====================================================================
\newif\ifswwb  \swwbfalse
\newif\ifswweb \swwebfalse
\makeatletter
% The workbook and website builds add a head above every question, so a page
% can no longer hold exactly what the paper's page holds.  In those builds a
% page may break only before a question head or, in a structured paper,
% before a part such as (b) or (ii) - never inside a question, between a stem
% and its options, or between a part and its answer lines.  The yearly build
% is untouched: there every \swnb is empty and every page break is the paper's.
\newcommand{\swnb}{\ifswwb\nobreak\interlinepenalty=10000 \else
  \ifswweb\nobreak\interlinepenalty=10000 \fi\fi}
\newcommand{\partbreak}{\par\ifswwb\penalty\z@\else\ifswweb\penalty\z@\fi\fi}
% the paper's page break before a question; in the head builds questions
% simply follow on, each kept whole
\newcommand{\swpage}{\par\newpage\pagestart}
\makeatother
\makeatletter
\newcommand{\swsep}{\kern4.2bp\textendash\kern4.2bp}
\gdef\swpaper{}\gdef\swqnum{}\gdef\swtopics{}\gdef\swmarks{}
\newcommand{\examq}[4]{\par % ends the paragraph before, so no stray space is left in it
  \gdef\swpaper{#1}\gdef\swqnum{#2}\gdef\swtopics{#3}\gdef\swmarks{#4}%
  \ifswwb\sw@qhead{\sw@wbhead{#1}{#2}{#3}}\else
  \ifswweb\sw@qhead{\sw@webhead{#1}{#2}{#3}{#4}}\fi\fi}
\newcommand{\sw@wbhead}[3]{{\bfseries #1\swsep Q#2}\hspace{8bp}%
  \raisebox{0.6bp}{\topicchips{#3}}}
\newcommand{\sw@webhead}[4]{{\bfseries #1\swsep Q#2\swsep #3\swsep #4}}
%  The head never takes space from the page, so every page of the workbook and
%  website builds holds exactly what the paper's page holds.  A question that
%  opens a page has its head set in the header band, under the logo (passed
%  there as an e-TeX mark).  A question that starts part-way down a page has
%  its head set in the gap the paper already leaves above it, 16bp above the
%  question's first line, so the question stays where the paper puts it.
\newif\ifsw@pending \sw@pendingfalse
\newdimen\sw@hd
\newcommand{\sw@qhead}[1]{%
  \ifdim\prevdepth<-999pt \marks5{\unexpanded{#1}}%
  \else \gdef\sw@pending{#1}\global\sw@pendingtrue \fi}
%  #1 = the measured distance to the next element, #2 = the room it keeps below the head.
%  The head's baseline goes #2 above the next element; when the paper's gap is
%  too tight for that (under 29bp), the head sits in the middle of the gap
%  instead, never closer than 13bp to the line above when there is room for it.
%  Either way the next element lands exactly where the paper puts it.
\newdimen\sw@rest
\def\sw@puthead#1#2{\global\sw@pendingfalse
  \sw@hd=\dimexpr #1bp-#2bp\relax
  \ifdim\sw@hd<13bp
    \sw@rest=\dimexpr #1bp/2\relax
    \ifdim\sw@rest>13bp \sw@rest=13bp\fi
    \ifdim\sw@hd<\sw@rest \sw@hd=\sw@rest\fi
  \fi
  \global\sw@rest=\dimexpr #1bp-\sw@hd\relax
  \vskip\dimexpr\sw@hd-\prevdepth-\ht\strutbox\relax
  \begingroup\leftskip\z@ \rightskip\z@ plus 1fil\@totalleftmargin\z@ \hangindent\z@
  % the chip is taller than a line of text; it must not push the question down
  \nointerlineskip\noindent\strut\raisebox{\z@}[\ht\strutbox][\dp\strutbox]{\sw@fit{\sw@pending}}%
  \par\endgroup}
%  one line: a head too long for the measure is scaled down to fit it
\newcommand{\sw@fit}[1]{\setbox\z@\hbox{\fontsize{11bp}{12bp}\selectfont #1}%
  \ifdim\wd\z@>\textwidth \resizebox{\textwidth}{!}{\box\z@}\else\box\z@\fi}
% the mark total for the question, read from \examq so there is one number
\makeatother

% ---- the question body ----------------------------------------------
%  Each element carries two numbers measured off the official paper: how far
%  its baseline sits below the one above it, and the column it starts in.
%  Blank working space needs no macro of its own - it is simply a larger
%  distance to the next element.
\makeatletter
\gdef\sw@hang{0}
%  a paragraph whose later lines line up under a phrase in the first
\newcommand{\hang}[1]{\par\gdef\sw@hang{#1}}
\newcommand{\sw@place}[3]{\par\swnb
  \leftskip=\dimexpr #2bp-49.6bp\relax
  \rightskip=\dimexpr 548.9bp-#3bp\relax
  \@totalleftmargin=\leftskip
  \linewidth=\dimexpr\textwidth-\leftskip-\rightskip\relax
  \ifsw@pending \sw@puthead{#1}{16}\nobreak
    \vskip\dimexpr\sw@rest-\prevdepth-\ht\strutbox\relax
  \else
  \ifdim\prevdepth<-999pt \vspace*{\dimexpr #1bp-\ht\strutbox\relax}
  \else \vskip\dimexpr #1bp-\prevdepth-\ht\strutbox\relax\fi\fi
  \ifdim\sw@hang bp>0pt \hangindent=\sw@hang bp\relax\hangafter=1\relax\fi
  \gdef\sw@hang{0}%
  \nointerlineskip\noindent\strut\ignorespaces}
% \prevdepth carries across \newpage, so the first element of a page would be
% spaced from the last line of the page before.  This resets it, which also
% tells \sw@place to set the distance from the top of the text area instead.
\newcommand{\pagestart}{\par\prevdepth=-1000pt\relax\marks5{}}
\newcommand{\swat}[1]{\sw@place{#1}{\sw@curx}{545.7}}
\def\sw@curx{72.3}
% a line of text starting in column #2
% no \par at the end: the next element's \sw@place closes the paragraph,
% so a line held back by \linebreak stays part of it
\newcommand{\row}[3]{\sw@place{#1}{#2}{545.7}#3}
% a label hanging #1 bp to the left of the text it belongs to
\newcommand{\lab}[2]{\llap{\makebox[#1bp][l]{\strut #2}}\ignorespaces}
% the paper's dotted answer line: 2.875bp between dots, as measured
\def\swdots{\leavevmode\cleaders\hb@xt@ 2.875bp{\hss.\hss}\hfill\kern\z@}
%   \dotline{gap}{from}{before the dots}{after them}{to}
%   These macros print nothing, so each one starts with \par: otherwise the
%   line break after it in the source would leave a space at the end of the
%   paragraph before, and a full line would be squeezed to make room for it.
%   \dotgrid{first dot}{last dot}{pitch}, given just before a \dotline, sets
%   its dots where the paper prints them: on the paper's own grid, stopping at
%   the last whole dot before the unit.  The unit then ends 5.5bp before the
%   mark, and the mark ends at the right margin.
\newif\ifsw@grid
\newcommand{\dotgrid}[3]{\par\gdef\sw@gfirst{#1}\gdef\sw@glast{#2}\gdef\sw@gpitch{#3}%
  \global\sw@gridtrue}
\def\sw@post#1[#2]#3\sw@nil{%
  \ifx\relax#1\relax\else{\ignorespaces#1\unskip}\fi
  \ifx\relax#2\relax\else\kern5.5bp[#2]\fi}
\newcommand{\dotline}[5]{\sw@place{#1}{#2}{#5}%
  \ifsw@grid\global\sw@gridfalse
    \rlap{\kern\dimexpr\sw@gfirst bp-#2bp\relax
      \hbox to\dimexpr\sw@glast bp-\sw@gfirst bp+\sw@gpitch bp+0.5bp\relax
        {\leaders\hbox to\sw@gpitch bp{.\hss}\hfill}}%
    #3\hfill\sw@post#4[]\sw@nil
  \else
    \ifx\relax#3\relax\else#3\kern3.4bp\fi\swdots
    \ifx\relax#4\relax\else\kern3.4bp#4\fi
  \fi}
% the mark for a part, set where the paper sets it.  Marks live only inside
% these macros, never as loose text, so the site can count them.
\newcommand{\swmark}[1]{\unskip\nobreak\hfill\mbox{#1}}
\newcommand{\qmark}[2]{\sw@place{#1}{49.6}{545.7}\hfill\mbox{#2}}
% the question total, taken from \examq so the number is written once
\newcommand{\total}[1]{\sw@place{#1}{49.6}{545.7}\hfill\mbox{[T\kern-1.0bp otal:\ \swmarks]}}% the paper kerns "To"
\newcommand{\gap}[1]{\par\vspace{#1bp}}
\makeatother

% ---- mathematics, placed piece by piece -----------------------------
%  A line of mathematics - fractions, radicals, symbols set apart - is set as
%  a \freerow: every run of letters is placed where the paper prints it,
%  measured in bp from the column at 49.6bp and from the row's baseline (up is
%  positive); fraction bars and the strokes of a radical sign are drawn where
%  the paper draws them; and a glyph no free font has (a stretched bracket,
%  MathType's own symbols) is cut out of the official PDF at its exact place.
\makeatletter
\newcommand{\freerow}[2]{\sw@place{#1}{49.6}{545.7}#2}
\newcommand{\at}[3]{\rlap{\kern\dimexpr #1bp-49.6bp\relax\raisebox{#2bp}[0pt][0pt]{#3}}}
\newcommand{\ruleat}[4]{\rlap{\kern\dimexpr #1bp-49.6bp\relax
  \raisebox{\dimexpr #3bp-#4bp/2\relax}[0pt][0pt]{\rule{\dimexpr #2bp-#1bp\relax}{#4bp}}}}
%  \swpath{style}{path}{clip}: a drawing, drawn through the paper's own
%  clipping path when it has one (a radical sign is a thick stroke clipped to
%  the sign's outline)
\newcommand{\swpath}[3]{\rlap{\begin{tikzpicture}[overlay,x=1bp,y=1bp,baseline=0pt]
  \begin{scope}#3\path[#1] #2;\end{scope}\end{tikzpicture}}}
\newcommand{\swglyph}[2]{\swclip{#1}{#2}}
% how far a band of mathematics hangs below its first baseline
\newcommand{\swdepth}[1]{\vrule width\z@ height\z@ depth#1bp\relax}
\makeatother
% a glyph from the real Symbol font, by its code: \swsy{70} is pi
\newcommand{\swsy}[1]{{\usefont{U}{psy}{m}{n}\char"#1\relax}}

% ---- figures, clipped out of the official PDF at compile time --------
%   \qfig{page}{l b r t}{caption}{above}{below}{left}
%   above = bp from the previous baseline down to the top of the figure
%   left  = the figure's left edge on the official page
\newcommand{\srcpdf}{4037_s16_qp_22.pdf}
\newcommand{\swclip}[3][]{%
  \IfFileExists{\srcpdf}{\includegraphics[page=#2,trim=#3,clip,#1]{\srcpdf}}%
  {\fbox{\parbox[c][22mm][c]{0.45\textwidth}{\centering\ttfamily\footnotesize
    source PDF missing:\\\detokenize\expandafter{\srcpdf}}}}}
\makeatletter
\newcommand{\qfig}[7][]{\par\swnb\begingroup\setlength{\parskip}{0pt}%
  \ifsw@pending \sw@puthead{#5}{12}\nobreak
    \vskip\dimexpr\sw@rest-\prevdepth\relax
  \else
  \ifdim\prevdepth<-999pt \vspace*{#5bp}
  \else \vskip\dimexpr #5bp-\prevdepth\relax\fi\fi
  \nointerlineskip
  \noindent\makebox[\linewidth][l]{\hspace*{-\@totalleftmargin}%
    \hspace*{\dimexpr #7bp-49.6bp\relax}\swclip[#1]{#2}{#3}\hfil}\par
  \ifx\relax#4\relax\else
    \vskip #6bp\relax\nointerlineskip
    \noindent\makebox[\linewidth][l]{\hspace*{-\@totalleftmargin}%
      \makebox[\textwidth][c]{\textbf{#4}}}\par
  \fi\endgroup}
\makeatother

% ---- the School Wits head -------------------------------------------
%  The official paper leaves the top of the page empty apart from a centred
%  page number, which we move to the foot. The head is built inside that
%  margin, so the question area below it never moves.
\newlength{\swlogoht}\setlength{\swlogoht}{14bp}
\newlength{\swheadgap}\setlength{\swheadgap}{3bp}
\newlength{\swqlinenow}% the line actually reserved: none in workbook mode,
                       % so the logo drops to sit just above the question head
\newlength{\swqline}\setlength{\swqline}{25.4bp}
% the head of a question that opens a page sits 16bp above the question's
% first line, the same as a head set in the gap above a question mid-page
\newlength{\swheaddrop}\setlength{\swheaddrop}{1.05bp}
\newcommand{\swlogofile}{3-01.png}
%  mark then hairline. For the mark on the right instead, use:
%    \renewcommand{\swlogoalign}{\hfil\swlogo}
\newcommand{\swlogoalign}{\swlogo\hspace{13bp}%
  \color{swrule}\leaders\hbox{\vrule height0.52\swlogoht depth-0.48\swlogoht width1bp}\hfil}
\newcommand{\swlogo}{%
  \IfFileExists{\swlogofile}%
    {\includegraphics[height=\swlogoht]{\swlogofile}}%
    {{\usefont{T1}{ptm}{b}{sc}\fontsize{11.5bp}{11.5bp}\selectfont School Wits}}}
\gdef\swrunhead{}
\makeatletter
% The logo band always sits in the top margin at the same place.  The line
% under it is optional: with \swrunhead empty the space is simply left blank,
% so the logo and the body below it never move.
\newcommand{\swheadblock}{%
  % the head of a question that opens the page, pinned to the header's
  % baseline: its own baseline lands 16bp above the question's first line
  \rlap{\raisebox{-\swheaddrop}[\z@][\z@]{\sw@fit{\botmarks5}}}%
  \parbox[b][\dimexpr\swlogoht+\swheadgap+\swqline\relax][t]{\textwidth}{%
    \hb@xt@\textwidth{\swlogoalign}}}
\makeatother
\pagestyle{fancy}\fancyhf{}
\renewcommand{\headrulewidth}{0pt}\renewcommand{\footrulewidth}{0pt}
\fancyhead[L]{\swheadblock}
\fancyfoot[C]{\ftsize\thepage}

\begin{document}
\normalsize
% ------------------------------------------------------- HEADER MODE --
\swwebtrue     % website build: this is the WEB file
% ----------------------------------------------------------------------
\gdef\swrunhead{}
\pagestart
\freerow{-0.18}{%
  \at{239.65}{0.00}{{\bfseries\itshape Mathematical} {\bfseries\itshape Formulae}}}
\partbreak
\freerow{78.0}{%
  \at{262.65}{-0.00}{\textbf{1.}}%
  \at{276.40}{-0.00}{\textbf{ALGEBRA}}}
\row{26.0}{49.62}{{\itshape Quadratic} {\itshape Equation}}
\row{26.01}{72.3}{For the equation {\itshape ax}\raisebox{4.04bp}[0pt][0pt]{\fontsize{8.62bp}{8.62bp}\selectfont 2} + {\itshape bx} + {\itshape c} = 0,}
\qfig{2}{252.30bp 604.90bp 251.22bp 205.33bp}{}{2.9}{0}{252.30}
\row{37.56}{49.61}{{\itshape Binomial} {\itshape Theor}\kern-0.42bp {\itshape em}}
\freerow{18.48}{%
  \at{218.41}{-19.30}{\swglyph{2}{218.41bp 529.56bp 373.67bp 284.65bp}}%
  \at{222.10}{-12.99}{1}%
  \at{229.68}{-19.30}{\swglyph{2}{229.68bp 529.56bp 362.39bp 284.65bp}}%
  \at{272.47}{-19.30}{\swglyph{2}{272.47bp 529.56bp 319.61bp 284.65bp}}%
  \at{276.55}{-12.99}{2}%
  \at{283.74}{-19.30}{\swglyph{2}{283.74bp 529.56bp 308.33bp 284.65bp}}%
  \at{354.34}{-19.30}{\swglyph{2}{354.34bp 529.56bp 237.74bp 284.65bp}}%
  \at{359.29}{-12.99}{{\itshape r}}%
  \at{365.61}{-19.30}{\swglyph{2}{365.61bp 529.56bp 226.46bp 284.65bp}}%
  \at{222.10}{0.00}{{\itshape n}}%
  \at{276.55}{0.00}{{\itshape n}}%
  \at{358.66}{0.00}{{\itshape n}}%
  \at{147.86}{-7.52}{({\itshape a} + {\itshape b})}%
  \at{179.26}{-3.49}{{\fontsize{8.62bp}{8.62bp}\selectfont {\itshape n}}}%
  \at{186.45}{-7.52}{= {\itshape a}}%
  \at{201.56}{-3.49}{{\fontsize{8.62bp}{8.62bp}\selectfont {\itshape n}}}%
  \at{208.74}{-7.52}{+}%
  \at{232.62}{-7.52}{{\itshape a}}%
  \at{238.37}{-3.49}{{\fontsize{8.62bp}{8.62bp}\selectfont {\itshape n}\textendash{}1}}%
  \at{254.18}{-7.52}{{\itshape b} +}%
  \at{286.68}{-7.52}{{\itshape a}}%
  \at{292.43}{-3.49}{{\fontsize{8.62bp}{8.62bp}\selectfont {\itshape n}\textendash{}2}}%
  \at{308.24}{-7.52}{{\itshape b}}%
  \at{313.99}{-3.49}{{\fontsize{8.62bp}{8.62bp}\selectfont 2}}%
  \at{321.18}{-7.52}{+ \swsy{BC} +}%
  \at{368.79}{-7.52}{{\itshape a}}%
  \at{374.54}{-3.49}{{\fontsize{8.62bp}{8.62bp}\selectfont {\itshape n}\textendash{}{\itshape r}}}%
  \at{389.40}{-7.52}{{\itshape b}}%
  \at{395.15}{-3.49}{{\fontsize{8.62bp}{8.62bp}\selectfont {\itshape r}}}%
  \at{401.38}{-7.52}{+ \swsy{BC} + {\itshape b}}%
  \at{440.22}{-3.49}{{\fontsize{8.62bp}{8.62bp}\selectfont {\itshape n}}}%
  \at{444.54}{-7.52}{,}%
  \swdepth{21.73}}
\partbreak
\freerow{39.0}{%
  \at{224.03}{-19.30}{\swglyph{2}{224.03bp 490.56bp 368.04bp 323.65bp}}%
  \at{228.99}{-12.99}{{\itshape r}}%
  \at{235.31}{-19.30}{\swglyph{2}{235.31bp 490.56bp 356.76bp 323.65bp}}%
  \at{228.36}{0.00}{{\itshape n}}%
  \at{266.73}{-1.44}{{\itshape n}!}%
  \at{72.28}{-7.51}{where {\itshape n} is a positive integer and}%
  \at{241.36}{-7.52}{=}%
  \at{250.76}{-14.44}{({\itshape n} \textendash{} {\itshape r})!{\itshape r}!}%
  \swpath{draw=black,line width=0.50bp}{(201.12,-4.53)--(242.72,-4.53)}{}%
  \swdepth{21.73}}
\partbreak
\freerow{85.51}{%
  \at{241.77}{0.00}{\textbf{2.}}%
  \at{255.94}{0.00}{\textbf{TRIGONOMETR}\kern-0.40bp \textbf{Y}}}
\row{26.0}{49.6}{{\itshape Identities}}
\freerow{26.0}{%
  \at{253.93}{-0.00}{sin}%
  \at{267.36}{4.02}{{\fontsize{8.62bp}{8.62bp}\selectfont 2}}%
  \at{274.55}{-0.01}{{\itshape A} + cos}%
  \at{309.14}{4.02}{{\fontsize{8.62bp}{8.62bp}\selectfont 2}}%
  \at{316.33}{-0.01}{{\itshape A} = 1}}
\freerow{21.0}{%
  \at{253.95}{-0.00}{sec}%
  \at{268.63}{4.02}{{\fontsize{8.62bp}{8.62bp}\selectfont 2}}%
  \at{275.82}{-0.00}{{\itshape A} = 1 + tan}%
  \at{327.11}{4.02}{{\fontsize{8.62bp}{8.62bp}\selectfont 2}}%
  \at{334.30}{-0.00}{{\itshape A}}}
\freerow{21.0}{%
  \at{248.52}{-0.00}{cosec}%
  \at{274.06}{4.02}{{\fontsize{8.62bp}{8.62bp}\selectfont 2}}%
  \at{281.25}{-0.00}{{\itshape A} = 1 + cot}%
  \at{332.54}{4.02}{{\fontsize{8.62bp}{8.62bp}\selectfont 2}}%
  \at{339.73}{-0.00}{{\itshape A}}}
\row{39.0}{49.61}{{\itshape Formulae} {\itshape for} \swsy{44}{\itshape ABC}}
\freerow{18.63}{%
  \at{257.09}{-0.00}{{\itshape a}}%
  \at{294.60}{-0.00}{{\itshape b}}%
  \at{332.60}{-0.00}{{\itshape c}}%
  \at{275.36}{-7.37}{=}%
  \at{313.11}{-7.37}{=}%
  \at{248.41}{-13.00}{sin {\itshape A}}%
  \at{285.82}{-13.00}{sin {\itshape B}}%
  \at{323.17}{-13.00}{sin {\itshape C}}%
  \swpath{draw=black,line width=0.50bp}{(197.84,-3.65)--(222.89,-3.65)}{}%
  \swpath{draw=black,line width=0.50bp}{(235.12,-3.65)--(260.64,-3.65)}{}%
  \swpath{draw=black,line width=0.50bp}{(272.87,-3.65)--(298.23,-3.65)}{}%
  \swdepth{15.30}}
\freerow{46.37}{%
  \at{242.53}{-0.00}{{\itshape a}}%
  \at{248.28}{4.02}{{\fontsize{8.62bp}{8.62bp}\selectfont 2}}%
  \at{255.46}{-0.00}{= {\itshape b}}%
  \at{270.58}{4.02}{{\fontsize{8.62bp}{8.62bp}\selectfont 2}}%
  \at{277.76}{-0.00}{+ {\itshape c}}%
  \at{292.23}{4.02}{{\fontsize{8.62bp}{8.62bp}\selectfont 2}}%
  \at{299.42}{-0.00}{\textendash{} 2{\itshape bc} cos {\itshape A}}}
\partbreak
\freerow{18.95}{%
  \at{284.54}{0.00}{1}%
  \at{264.17}{-7.06}{\swsy{44} =}%
  \at{294.27}{-7.05}{{\itshape bc} sin {\itshape A}}%
  \at{284.54}{-12.99}{2}%
  \swpath{draw=black,line width=0.50bp}{(234.84,-3.08)--(240.79,-3.08)}{}%
  \swdepth{15.29}}
\newpage
\pagestart
\examq{4037/22/M/J/16}{1}{Quadratic Functions}{4}
\freerow{3.69}{%
  \at{49.61}{0.00}{\textbf{1}}%
  \at{72.77}{0.00}{\textbf{(i)}}%
  \at{94.96}{0.00}{Given that\kern7.25bp {\itshape x}}%
  \at{156.09}{5.16}{{\fontsize{8.00bp}{8.00bp}\selectfont 2}}%
  \at{162.47}{0.00}{{\itshape +}\kern2.37bp 2{\itshape kx}\kern2.35bp {\itshape +}\kern2.37bp 4{\itshape k}}%
  \at{214.86}{0.48}{\swsy{2D}}%
  \at{224.56}{0.00}{3\kern3.15bp {\itshape =}\kern3.15bp 0\kern6.44bp has no real roots, show that {\itshape k} satisfies\kern6.73bp {\itshape k}}%
  \at{442.74}{5.16}{{\fontsize{8.00bp}{8.00bp}\selectfont 2}}%
  \at{449.59}{0.48}{\swsy{2D}}%
  \at{459.30}{0.00}{4{\itshape k}\kern2.36bp {\itshape +}\kern2.37bp 3}%
  \at{492.47}{-0.22}{\swglyph{3}{492.47bp 765.38bp 94.56bp 68.03bp}}%
  \at{504.30}{0.00}{0\kern0.69bp .}%
  \at{532.27}{0.00}{\swmark{[2]}}}
\partbreak
\freerow{155.99}{%
  \at{69.57}{-0.00}{\textbf{(ii)}}%
  \at{94.97}{-0.00}{Solve the inequality}%
  \at{214.70}{-0.01}{{\itshape k}}%
  \at{220.36}{5.15}{{\fontsize{8.00bp}{8.00bp}\selectfont 2}}%
  \at{227.21}{0.47}{\swsy{2D}}%
  \at{236.92}{-0.01}{4{\itshape k}\kern2.36bp {\itshape +}\kern2.37bp 3}%
  \at{270.09}{-0.23}{\swglyph{3}{270.09bp 609.38bp 317.18bp 224.27bp}}%
  \at{281.92}{-0.01}{0\kern0.69bp .}%
  \at{532.26}{-0.01}{\swmark{[2]}}}
\examq{4037/22/M/J/16}{2}{Differentiation}{5}
\freerow{162.97}{%
  \at{302.92}{-0.00}{5{\itshape x}}%
  \at{316.61}{0.47}{\swsy{2D}}%
  \at{326.32}{-0.00}{1}%
  \at{49.60}{-6.04}{\textbf{2}}%
  \at{72.28}{-6.04}{V\kern-1.28bp ariables {\itshape x} and {\itshape y} are related by the equation\kern7.79bp {\itshape y}\kern3.14bp {\itshape =}}%
  \at{335.17}{-6.04}{.}%
  \at{305.79}{-12.29}{3}%
  \at{314.40}{-11.81}{\swsy{2D}}%
  \at{324.10}{-12.29}{{\itshape x}}%
  \at{120.58}{-24.13}{d{\itshape y}}%
  \at{72.77}{-32.04}{\textbf{(i)}}%
  \at{94.96}{-32.04}{Find}%
  \at{133.08}{-32.04}{, simplifying your answer\kern-0.64bp .}%
  \at{120.58}{-38.50}{d{\itshape x}}%
  \ruleat{302.35}{332.63}{-2.57}{0.50}%
  \ruleat{120.02}{131.99}{-28.57}{0.50}%
  \at{532.26}{-32.04}{\swmark{[2]}}%
  \swdepth{40.80}}
\partbreak
\freerow{175.04}{%
  \at{69.57}{0.00}{\textbf{(ii)}}%
  \at{94.96}{0.00}{Hence find the approximate change in {\itshape x} when {\itshape y} increases from 9 by the small amount 0.07.}%
  \at{532.26}{0.00}{\swmark{[3]}}}
\newpage
\pagestart
\examq{4037/22/M/J/16}{3}{Permutations \& Combinations}{3}
\row{0.0}{72.28}{\lab{22.67}{\textbf{3}}A\kern2.34bp team of 3 people is to be selected from 7 women and 6 men. Find the number of dif\kern-0.17bp ferent teams that\linebreak}
could be selected if there must be more women than men on the team.\swmark{[3]}
\examq{4037/22/M/J/16}{4}{Factors of Polynomials}{5}
\row{195.0}{72.28}{\lab{22.67}{\textbf{4}}\textbf{Do} \textbf{not} \textbf{use} \textbf{a} \textbf{calculator} \textbf{in} \textbf{this} \textbf{question.}}
\freerow{26.0}{%
  \at{72.28}{-0.01}{The polynomial\kern4.00bp p\kern1.18bp ({\itshape x})\kern3.18bp {\itshape =}\kern3.15bp 2{\itshape x}}%
  \at{194.67}{5.15}{{\fontsize{8.00bp}{8.00bp}\selectfont 3}}%
  \at{201.53}{0.48}{\swsy{2D}}%
  \at{211.23}{-0.00}{3{\itshape x}}%
  \at{222.64}{5.15}{{\fontsize{8.00bp}{8.00bp}\selectfont 2}}%
  \at{229.01}{-0.00}{{\itshape +}\kern2.37bp {\itshape qx}\kern2.35bp {\itshape +}\kern2.37bp 56\kern3.53bp has a factor {\itshape x} \textendash{} 2.}}
\partbreak
\row{26.01}{94.96}{\lab{22.2}{\textbf{(i)}}Show that {\itshape q} = \textminus{}30.\swmark{[1]}}
\partbreak
\freerow{116.99}{%
  \at{69.57}{-0.01}{\textbf{(ii)}}%
  \at{94.96}{-0.01}{Factorise\kern4.00bp p\kern1.18bp ({\itshape x})\kern3.55bp completely and hence state all the solutions of\kern3.99bp p\kern1.18bp ({\itshape x})\kern3.18bp {\itshape =}\kern3.14bp 0\kern3.58bp .}%
  \at{532.27}{-0.00}{\swmark{[4]}}}
\newpage
\pagestart
\examq{4037/22/M/J/16}{5}{Straight-Line Graphs}{7}
\row{0.02}{72.28}{\lab{22.67}{\textbf{5}}The coordinates of three points are {\itshape A}(\textminus{}2, 6), {\itshape B}(6, 10) and {\itshape C}({\itshape p}, 0).}
\partbreak
\row{26.0}{94.96}{\lab{22.19}{\textbf{(i)}}Find the coordinates of {\itshape M}, the mid-point of {\itshape AB}.\swmark{[2]}}
\partbreak
\freerow{129.99}{%
  \at{69.57}{-0.00}{\textbf{(ii)}}%
  \at{94.96}{-0.00}{Given that {\itshape CM} is perpendicular to {\itshape AB}, find the value of the constant {\itshape p}.}%
  \at{532.26}{-0.00}{\swmark{[2]}}}
\partbreak
\freerow{195.01}{%
  \at{66.37}{0.00}{\textbf{(iii)}}%
  \at{94.95}{0.00}{Find angle {\itshape MCB}.}%
  \at{532.26}{0.00}{\swmark{[3]}}}
\newpage
\pagestart
\examq{4037/22/M/J/16}{6}{Circular Measure}{7}
\row{-0.18}{49.61}{\textbf{6}}
\qfig{6}{226.74bp 652.83bp 228.90bp 78.22bp}{}{5.8}{0}{226.74}
\row{27.19}{72.28}{The diagram shows a sector of a circle with centre {\itshape O} and radius 5\kern1.48bp cm. The length of the arc {\itshape AB} is 7\kern1.47bp cm.}
\freerow{13.0}{%
  \at{72.30}{0.00}{Angle {\itshape AOB} is}%
  \at{139.43}{-0.70}{\swglyph{6}{139.43bp 611.94bp 449.76bp 220.99bp}}%
  \at{147.57}{0.00}{radians.}}
\partbreak
\freerow{26.0}{%
  \at{72.78}{-0.00}{\textbf{(i)}}%
  \at{94.97}{-0.00}{Explain why}%
  \at{156.68}{-0.94}{\swglyph{6}{156.68bp 585.70bp 432.51bp 247.23bp}}%
  \at{164.82}{-0.00}{must be greater than 1 radian.}%
  \at{532.28}{-0.00}{\swmark{[1]}}}
\partbreak
\freerow{91.0}{%
  \at{69.59}{-0.00}{\textbf{(ii)}}%
  \at{94.99}{-0.00}{Find the value of}%
  \at{176.25}{-0.94}{\swglyph{6}{176.25bp 494.70bp 412.95bp 338.23bp}}%
  \at{181.75}{-0.00}{.}%
  \at{532.28}{-0.00}{\swmark{[2]}}}
\partbreak
\freerow{143.0}{%
  \at{66.40}{-0.00}{\textbf{(iii)}}%
  \at{94.97}{-0.00}{Calculate the area of the sector\kern2.86bp {\itshape AOB}.}%
  \at{532.27}{-0.00}{\swmark{[2]}}}
\partbreak
\freerow{156.0}{%
  \at{67.03}{-0.00}{\textbf{(iv)}}%
  \at{94.97}{-0.00}{Calculate the area of the shaded segment.}%
  \at{532.27}{-0.00}{\swmark{[2]}}}
\newpage
\pagestart
\examq{4037/22/M/J/16}{7}{Excluded from the present curriculum}{8}
\freerow{-0.84}{%
  \at{148.41}{-23.02}{\swglyph{7}{148.41bp 747.11bp 443.18bp 65.42bp}}%
  \at{175.56}{-23.02}{\swglyph{7}{175.56bp 747.11bp 416.04bp 65.42bp}}%
  \at{273.83}{-23.02}{\swglyph{7}{273.83bp 747.11bp 317.52bp 65.42bp}}%
  \at{300.98}{-23.02}{\swglyph{7}{300.98bp 747.11bp 290.62bp 65.42bp}}%
  \at{152.22}{-4.53}{4}%
  \at{169.47}{-4.53}{5}%
  \at{277.64}{-4.53}{4}%
  \at{294.89}{-4.53}{2}%
  \at{148.41}{-23.04}{\swglyph{7}{148.41bp 747.09bp 443.18bp 65.44bp}}%
  \at{175.56}{-23.04}{\swglyph{7}{175.56bp 747.09bp 416.04bp 65.44bp}}%
  \at{273.83}{-23.04}{\swglyph{7}{273.83bp 747.09bp 317.52bp 65.44bp}}%
  \at{300.98}{-23.04}{\swglyph{7}{300.98bp 747.09bp 290.62bp 65.44bp}}%
  \at{49.61}{-12.37}{\textbf{7}}%
  \at{72.28}{-12.37}{The matrix \textbf{A} is}%
  \at{148.41}{-23.16}{\swglyph{7}{148.41bp 746.97bp 443.18bp 65.56bp}}%
  \at{175.56}{-23.16}{\swglyph{7}{175.56bp 746.97bp 416.04bp 65.56bp}}%
  \at{182.69}{-12.36}{and the matrix \textbf{B} is}%
  \at{273.83}{-23.16}{\swglyph{7}{273.83bp 746.97bp 317.52bp 65.56bp}}%
  \at{300.98}{-23.16}{\swglyph{7}{300.98bp 746.97bp 290.62bp 65.56bp}}%
  \at{305.23}{-12.36}{.}%
  \at{152.22}{-19.46}{3}%
  \at{169.47}{-19.46}{2}%
  \at{277.64}{-19.46}{1}%
  \at{294.89}{-19.46}{3}%
  \at{148.41}{-23.01}{\swglyph{7}{148.41bp 747.12bp 443.18bp 65.41bp}}%
  \at{175.56}{-23.01}{\swglyph{7}{175.56bp 747.12bp 416.04bp 65.41bp}}%
  \at{273.83}{-23.01}{\swglyph{7}{273.83bp 747.12bp 317.52bp 65.41bp}}%
  \at{300.98}{-23.01}{\swglyph{7}{300.98bp 747.12bp 290.62bp 65.41bp}}%
  \swdepth{27.01}}
\partbreak
\row{38.36}{94.97}{\lab{22.2}{\textbf{(i)}}Find the matrix \textbf{C} such that\kern3.54bp \textbf{C}\kern3.16bp {\itshape =}\kern3.15bp 3\textbf{A}\kern2.38bp {\itshape +}\kern2.37bp \textbf{B}\kern0.64bp .\swmark{[2]}}
\partbreak
\freerow{182.01}{%
  \at{69.57}{0.00}{\textbf{(ii)}}%
  \at{94.96}{0.00}{Show that\kern3.57bp det\kern1.16bp (\textbf{AB})\kern3.18bp {\itshape =}\kern3.16bp det\kern1.54bp \textbf{A}}%
  \at{224.60}{-0.00}{\swsy{B4}}%
  \at{233.32}{-0.00}{det\kern1.54bp \textbf{B}\kern0.64bp .}%
  \at{532.26}{-0.00}{\swmark{[4]}}}
\partbreak
\freerow{321.39}{%
  \at{66.38}{-0.00}{\textbf{(iii)}}%
  \at{94.95}{-0.00}{Find the matrix\kern3.54bp (\textbf{AB})}%
  \at{194.58}{5.49}{{\fontsize{8.00bp}{8.00bp}\selectfont \swsy{2D}}}%
  \at{199.68}{5.16}{{\fontsize{8.00bp}{8.00bp}\selectfont 1}}%
  \at{204.78}{0.00}{.}%
  \at{532.26}{0.00}{\swmark{[2]}}}
\newpage
\pagestart
\examq{4037/22/M/J/16}{8}{Simultaneous Equations}{6}
\partbreak
\freerow{0.91}{%
  \at{376.13}{0.00}{5}%
  \at{395.69}{0.00}{3}%
  \at{49.61}{-6.03}{\textbf{8}}%
  \at{72.28}{-6.03}{Find the coordinates of the points of intersection of the curve\kern3.80bp 4\kern2.37bp {\itshape +}}%
  \at{384.94}{-6.03}{{\itshape +}}%
  \at{405.16}{-6.03}{{\itshape =}\kern3.15bp 0\kern3.56bp and the line}%
  \at{376.91}{-11.93}{{\itshape y}}%
  \at{396.22}{-11.93}{{\itshape x}}%
  \at{74.32}{-19.03}{{\itshape y}\kern3.14bp {\itshape =}\kern3.15bp 15{\itshape x}\kern2.36bp {\itshape +}\kern2.36bp 10\kern0.69bp .}%
  \ruleat{375.41}{382.57}{-2.56}{0.50}%
  \ruleat{395.13}{402.01}{-2.56}{0.50}%
  \at{532.26}{-19.03}{\swmark{[6]}}%
  \swdepth{21.33}}
\newpage
\pagestart
\examq{4037/22/M/J/16}{9}{Integration}{7}
\freerow{3.69}{%
  \at{134.93}{0.00}{{\itshape x}}%
  \at{140.58}{5.16}{{\fontsize{8.00bp}{8.00bp}\selectfont 3}}%
  \at{146.96}{0.00}{{\itshape +}\kern2.37bp {\itshape x}}%
  \at{162.80}{5.16}{{\fontsize{8.00bp}{8.00bp}\selectfont 2}}%
  \at{169.18}{0.00}{{\itshape +}\kern2.37bp 1}%
  \at{49.61}{-6.03}{\textbf{9}}%
  \at{72.28}{-6.03}{\textbf{(a)}}%
  \at{94.96}{-6.03}{Find}%
  \at{125.59}{-13.60}{\swglyph{9}{125.59bp 752.01bp 461.93bp 68.21bp}}%
  \at{160.86}{-9.90}{{\fontsize{8.00bp}{8.00bp}\selectfont 2}}%
  \at{187.17}{-6.02}{d{\itshape x}\kern1.00bp .}%
  \at{155.21}{-15.06}{{\itshape x}}%
  \ruleat{133.96}{185.68}{-2.56}{0.50}%
  \at{532.26}{-6.03}{\swmark{[3]}}%
  \swdepth{17.36}}
\partbreak
\freerow{162.03}{%
  \at{72.28}{0.00}{\textbf{(b)}}%
  \at{95.44}{0.00}{\textbf{(i)}}%
  \at{117.63}{0.00}{Find}%
  \at{142.28}{-7.56}{\swglyph{9}{142.28bp 596.01bp 445.24bp 224.20bp}}%
  \at{150.88}{0.00}{sin\kern1.17bp (5{\itshape x}\kern2.37bp {\itshape +}}%
  \at{192.86}{-0.70}{\swglyph{9}{192.86bp 602.87bp 395.37bp 232.46bp}}%
  \at{199.23}{0.00}{)\kern1.07bp d{\itshape x}\kern1.09bp .}%
  \at{532.26}{0.00}{\swmark{[2]}}%
  \swdepth{4.36}}
\partbreak
\freerow{182.0}{%
  \at{200.19}{8.85}{{\fontsize{8.00bp}{8.00bp}\selectfont 0}}%
  \at{92.25}{0.00}{\textbf{(ii)}}%
  \at{117.64}{0.00}{Hence evaluate}%
  \at{191.97}{-7.32}{\swglyph{9}{191.97bp 414.25bp 395.54bp 405.96bp}}%
  \at{203.45}{-4.61}{\swglyph{9}{203.45bp 416.96bp 387.42bp 421.01bp}}%
  \at{208.87}{0.00}{sin\kern1.17bp (5{\itshape x}\kern2.38bp {\itshape +}}%
  \at{250.86}{-0.70}{\swglyph{9}{250.86bp 420.87bp 337.38bp 414.46bp}}%
  \at{257.22}{0.00}{)\kern1.20bp d{\itshape x}\kern1.09bp .}%
  \at{197.15}{-7.60}{{\fontsize{8.00bp}{8.00bp}\selectfont \swsy{2D}}}%
  \at{203.59}{-11.56}{{\fontsize{7.00bp}{7.00bp}\selectfont 5}}%
  \ruleat{202.94}{207.72}{-5.54}{0.41}%
  \at{532.26}{0.00}{\swmark{[2]}}%
  \swdepth{12.96}}
\newpage
\pagestart
\examq{4037/22/M/J/16}{10}{Logarithmic \& Exponential Functions \textperiodcentered\ Equations, Inequalities \& Graphs}{9}
\freerow{3.69}{%
  \at{49.61}{0.00}{\textbf{10}}%
  \at{72.28}{0.00}{\textbf{(a)}}%
  \at{94.96}{0.00}{The\kern3.09bp graph\kern3.08bp of\kern3.08bp the\kern3.09bp curve\kern5.12bp {\itshape y}\kern3.14bp {\itshape =}\kern3.15bp {\itshape p}\kern1.18bp (4}%
  \at{241.95}{5.16}{{\fontsize{8.00bp}{8.00bp}\selectfont 2{\itshape x}}}%
  \at{249.52}{0.00}{)}%
  \at{256.22}{0.48}{\swsy{2D}}%
  \at{265.93}{0.00}{{\itshape q}\kern1.18bp (4}%
  \at{283.30}{5.16}{{\fontsize{8.00bp}{8.00bp}\selectfont {\itshape x}}}%
  \at{286.86}{0.00}{)\kern3.75bp passes\kern3.08bp through\kern3.08bp the\kern3.09bp points\kern3.08bp (0,\kern3.08bp 2)\kern3.08bp and\kern3.08bp (0.5,\kern3.08bp 14).\kern3.08bp Find\kern3.08bp the}}
\row{12.99}{94.96}{value of {\itshape p} and of {\itshape q}.\swmark{[3]}}
\partbreak
\freerow{195.01}{%
  \at{72.29}{0.00}{\textbf{(b)}}%
  \at{94.97}{0.00}{The\kern4.30bp variables\kern4.31bp {\itshape x}\kern4.31bp and\kern4.30bp {\itshape y}\kern4.31bp are\kern4.30bp connected\kern4.31bp by\kern4.31bp the\kern4.30bp equation\kern6.34bp {\itshape y}\kern3.14bp {\itshape =}\kern3.15bp 10}%
  \at{383.36}{5.16}{{\fontsize{8.00bp}{8.00bp}\selectfont 2{\itshape x}}}%
  \at{393.78}{0.48}{\swsy{2D}}%
  \at{403.48}{0.00}{2\kern1.18bp (10}%
  \at{426.60}{5.16}{{\fontsize{8.00bp}{8.00bp}\selectfont {\itshape x}}}%
  \at{430.16}{0.00}{)\kern0.68bp .\kern4.31bp Using\kern4.31bp the\kern4.30bp substitution}}
\row{13.0}{95.74}{{\itshape u}\kern3.16bp {\itshape =}\kern3.15bp 10\kern0.85bp \raisebox{5.16bp}[0pt][0pt]{\fontsize{8.00bp}{8.00bp}\selectfont {\itshape x}}\kern1.10bp , or otherwise, find the exact value of {\itshape x} when {\itshape y} = 24.\swmark{[3]}}
\partbreak
\freerow{195.0}{%
  \at{72.28}{-0.00}{\textbf{(c)}}%
  \at{94.95}{-0.00}{Solve\kern9.38bp log}%
  \at{145.56}{-4.00}{{\fontsize{8.00bp}{8.00bp}\selectfont 2}}%
  \at{150.75}{0.00}{({\itshape x}\kern2.38bp {\itshape +}\kern2.36bp 1)}%
  \at{184.70}{0.48}{\swsy{2D}}%
  \at{194.40}{0.00}{log}%
  \at{209.43}{-4.00}{{\fontsize{8.00bp}{8.00bp}\selectfont 2}}%
  \at{214.59}{0.00}{{\itshape x}\kern3.15bp {\itshape =}\kern3.14bp 3\kern3.01bp .}%
  \at{532.26}{0.00}{\swmark{[3]}}%
  \swdepth{5.60}}
\newpage
\pagestart
\examq{4037/22/M/J/16}{11}{Functions \textperiodcentered\ Quadratic Functions}{11}
\row{0.02}{94.96}{\lab{45.35}{\textbf{1}\kern-0.63bp \textbf{1}}\lab{22.68}{\textbf{(a)}}A\kern2.24bp function f is defined, for all real {\itshape x}, by}
\freerow{26.0}{%
  \at{266.59}{0.00}{f\kern1.20bp ({\itshape x})\kern3.17bp {\itshape =}\kern3.15bp {\itshape x}}%
  \at{306.48}{0.48}{\swsy{2D}}%
  \at{316.19}{0.00}{{\itshape x}}%
  \at{321.84}{5.16}{{\fontsize{8.00bp}{8.00bp}\selectfont 2}}%
  \at{327.62}{-0.52}{\swglyph{11}{327.62bp 742.75bp 265.66bp 97.14bp}}}
\row{26.0}{94.96}{Find the greatest value of f\kern1.44bp ({\itshape x}) and the value of {\itshape x} for which this occurs.\swmark{[3]}}
\partbreak
\freerow{129.99}{%
  \at{72.28}{-0.00}{\textbf{(b)}}%
  \at{94.96}{-0.00}{The\kern3.66bp domain\kern3.66bp of\kern4.40bp g\kern1.18bp ({\itshape x})\kern3.18bp {\itshape =}\kern3.15bp {\itshape x}}%
  \at{210.43}{0.47}{\swsy{2D}}%
  \at{220.14}{-0.01}{{\itshape x}}%
  \at{225.80}{5.15}{{\fontsize{8.00bp}{8.00bp}\selectfont 2}}%
  \at{234.55}{-0.01}{is\kern3.66bp such\kern3.66bp that\kern4.42bp g}%
  \at{298.92}{5.49}{{\fontsize{8.00bp}{8.00bp}\selectfont \swsy{2D}}}%
  \at{304.02}{5.15}{{\fontsize{8.00bp}{8.00bp}\selectfont 1}}%
  \at{309.21}{-0.01}{({\itshape x})\kern4.33bp exists.\kern3.66bp Explain\kern3.65bp why\kern5.17bp {\itshape x}}%
  \at{431.82}{-0.71}{\swglyph{11}{431.82bp 586.57bp 155.70bp 246.60bp}}%
  \at{442.56}{-0.01}{1\kern3.35bp is\kern3.66bp a\kern3.66bp suitable\kern3.65bp domain}}
\row{13.0}{94.96}{for g({\itshape x}).\swmark{[1]}}
\partbreak
\row{103.99}{94.96}{\lab{22.68}{\textbf{(c)}}The functions h and k are defined by}
\freerow{13.01}{%
  \at{271.71}{-0.00}{h:\kern1.17bp {\itshape x}}%
  \at{290.23}{0.50}{\swglyph{11}{290.23bp 457.78bp 296.57bp 380.67bp}}%
  \at{301.86}{-0.00}{lg\kern1.18bp ({\itshape x}\kern2.38bp {\itshape +}\kern2.36bp 2)}%
  \at{375.59}{-0.00}{for\kern4.37bp {\itshape x}}%
  \at{402.48}{-0.47}{\swglyph{11}{402.48bp 456.81bp 184.56bp 376.60bp}}%
  \at{412.63}{0.47}{\swsy{2D}}%
  \at{420.96}{-0.01}{2\kern0.60bp ,}}
\freerow{18.67}{%
  \at{271.71}{-0.00}{k:\kern1.17bp {\itshape x}}%
  \at{290.22}{0.74}{\swglyph{11}{290.22bp 439.35bp 296.57bp 399.10bp}}%
  \at{301.86}{-0.00}{5\kern2.38bp {\itshape +}}%
  \at{328.55}{-0.00}{{\itshape x}}%
  \at{336.51}{0.48}{\swsy{2D}}%
  \at{346.21}{-0.00}{1}%
  \at{375.59}{-0.00}{for\kern2.65bp 1}%
  \at{401.41}{-0.22}{\swglyph{11}{401.41bp 438.39bp 185.86bp 395.26bp}}%
  \at{413.25}{-0.00}{{\itshape x}}%
  \at{422.36}{-0.22}{\swglyph{11}{422.36bp 438.39bp 164.92bp 395.26bp}}%
  \at{434.18}{-0.00}{101\kern-0.31bp .}%
  \swpath{draw=black,line width=1.80bp}{(271.28,4.04)--(272.12,4.53)--(274.76,-0.22)}{\clip (274.06,9.68)--(274.06,6.09)--(271.02,4.30)--(271.28,4.04)--(272.12,4.53)--(274.76,-0.22)--(277.89,9.18)--(302.36,9.18)--(302.36,9.68)--(274.06,9.68);}%
  \swpath{draw=black,line width=1.00bp}{(274.76,-0.22)--(277.89,9.18)--(302.36,9.18)}{\clip (274.06,9.68)--(274.06,6.09)--(271.02,4.30)--(271.28,4.04)--(272.12,4.53)--(274.76,-0.22)--(277.89,9.18)--(302.36,9.18)--(302.36,9.68)--(274.06,9.68);}}
\partbreak
\row{26.01}{117.64}{\lab{22.2}{\textbf{(i)}}Find hk(10).\swmark{[2]}}
\partbreak
\freerow{103.99}{%
  \at{92.24}{-0.00}{\textbf{(ii)}}%
  \at{117.64}{-0.00}{Find\kern3.88bp k}%
  \at{149.25}{5.49}{{\fontsize{8.00bp}{8.00bp}\selectfont \swsy{2D}}}%
  \at{154.35}{5.15}{{\fontsize{8.00bp}{8.00bp}\selectfont 1}}%
  \at{159.54}{-0.01}{({\itshape x})\kern0.67bp , stating its domain and range.}%
  \at{532.26}{-0.00}{\swmark{[5]}}}
\freerow{247.0}{%
  \at{200.52}{-0.00}{\textbf{Question} \textbf{12} \textbf{is} \textbf{printed} \textbf{on} \textbf{the} \textbf{next} \textbf{page.}}}
\newpage
\pagestart
\examq{4037/22/M/J/16}{12}{Trigonometry}{8}
\row{0.0}{72.28}{\lab{22.67}{\textbf{12}}Solve the equation}
\partbreak
\freerow{25.84}{%
  \at{72.77}{-0.17}{\textbf{(i)}}%
  \at{95.36}{-0.16}{8\kern1.56bp sin}%
  \at{116.65}{4.99}{{\fontsize{8.00bp}{8.00bp}\selectfont 2}}%
  \at{121.21}{-0.16}{{\itshape A}\kern2.39bp {\itshape +}\kern2.36bp 2\kern1.56bp cos\kern1.54bp {\itshape A}\kern3.17bp {\itshape =}\kern3.15bp 7\kern6.31bp for\kern6.40bp 0}%
  \at{224.31}{5.06}{\swglyph{12}{224.31bp 748.51bp 368.01bp 89.94bp}}%
  \at{230.20}{-0.86}{\swglyph{12}{230.20bp 742.59bp 357.32bp 90.58bp}}%
  \at{240.93}{-0.16}{{\itshape A}}%
  \at{251.26}{-0.86}{\swglyph{12}{251.26bp 742.59bp 336.25bp 90.58bp}}%
  \at{261.99}{-0.16}{180}%
  \at{279.77}{5.06}{\swglyph{12}{279.77bp 748.51bp 312.55bp 89.94bp}}%
  \at{283.46}{-0.16}{,}%
  \at{532.25}{-0.16}{\swmark{[4]}}}
\partbreak
\freerow{273.16}{%
  \at{69.56}{-0.00}{\textbf{(ii)}}%
  \at{95.66}{-0.00}{cosec\kern1.14bp (3{\itshape B}\kern2.40bp {\itshape +}\kern2.38bp 1)\kern3.16bp {\itshape =}\kern3.15bp 2.5\kern8.93bp for\kern9.29bp 0}%
  \at{230.29}{-0.70}{\swglyph{12}{230.29bp 469.59bp 357.47bp 363.58bp}}%
  \at{241.03}{-0.00}{{\itshape B}}%
  \at{251.36}{-0.70}{\swglyph{12}{251.36bp 469.59bp 336.16bp 363.58bp}}%
  \at{262.23}{-0.70}{\swglyph{12}{262.23bp 469.59bp 326.01bp 365.74bp}}%
  \at{272.78}{-0.00}{radians.}%
  \at{532.26}{-0.00}{\swmark{[4]}}}
\end{document}
